\sectiondivider{Parte II}{L'Interactive Theorem Prover Lean 4}

% Slide 01 (Lean4)
\begin{frame}{Calcolo delle Costruzioni Induttive (CiC)}
    \begin{itemize}
        \item Lean 4 è basato sul \textbf{Calculus of Inductive Constructions} (CiC).
        \item Gerarchia rigorosa degli universi per prevenire paradossi (es. Girard):
        \[
        \mathtt{Prop} \equiv \mathtt{Sort}\ 0 \quad\prec\quad \mathtt{Type}\ 0 \equiv \mathtt{Sort}\ 1 \quad\prec\quad \mathtt{Type}\ 1 \quad\dots
        \]
    \end{itemize}

    \vskip 0.20cm
    \begin{columns}[t, onlytextwidth]
        \begin{column}{0.48\textwidth}
            \begin{block}{Universo Prop}
                \begin{minipage}[t][2.5cm][t]{\linewidth}
                    \begin{itemize}
                        \item Riservato alle asserzioni logiche.
                        \item \textbf{Proof Irrelevance}: due dimostrazioni della stessa proposizione sono sintatticamente indistinguibili.
                        \item Cancellazione a runtime (zero overhead).
                    \end{itemize}
                \end{minipage}
            \end{block}
        \end{column}

        \begin{column}{0.48\textwidth}
            \begin{block}{Gerarchia Type}
                \begin{minipage}[t][2.5cm][t]{\linewidth}
                    \begin{itemize}
                        \item Riservata ai dati e alle computazioni.
                        \item Contenuto computazionale rilevante (es. numeri, liste, alberi di sintassi).
                        \item Eseguibile ed estraibile in codice nativo C.
                    \end{itemize}
                \end{minipage}
            \end{block}
        \end{column}
    \end{columns}
\end{frame}

% Slide 02 (Lean4)
\begin{frame}{Il legame tra CiC e Logica in Lean 4}
    \begin{itemize}
        \item Il \textbf{Calculus of Inductive Constructions} (CiC) non è solo un linguaggio di programmazione, ma la formulazione computazionale della \textbf{logica intuizionista di ordine superiore} (IHOL).
        \item In Lean 4 la logica non è un modulo aggiuntivo, ma emerge nativamente dal calcolo dei tipi:
    \end{itemize}

    \begin{block}{I due pilastri logici del CiC}
        \begin{itemize}
            \item \textbf{Tipi Dipendenti ($\Pi$)}: generalizzano l'implicazione alla quantificazione universale ($\forall x. P(x)$). Il passaggio logico dall'ipotesi alla tesi coincide con l'applicazione funzionale.
            \item \textbf{Tipi Induttivi}: definiscono i predicati e i connettivi logici. Le regole di introduzione sono i \emph{costruttori}, l'eliminazione coincide con l'induzione/ricorsione (\texttt{.rec}).
        \end{itemize}
    \end{block}

    \begin{block}{Conseguenza fondamentale}
        In Lean 4 \textbf{non esiste un motore di inferenza logica}: verificare la validità di una dimostrazione coincide interamente con il \emph{type checking} del kernel sul termine generato.
    \end{block}
\end{frame}

% Slide 03 (Lean4)
\begin{frame}{Architettura TCB e Principio di de Bruijn}
    \begin{block}{Small Trusted Kernel (TCB minimale)}
        L'affidabilità dell'intero sistema non dipende dalla complessità del compilatore o delle tattiche, ma esclusivamente da un kernel compatto.
    \end{block}

    \vskip 0.20cm
    \begin{itemize}
        \item \textbf{Principio di de Bruijn}:
        \begin{itemize}
            \item Le tattiche (automazione, metaprogrammazione) generano $\lambda$-termini espliciti.
            \item Solo il micro-kernel verifica la tipabilità del termine finale prima di certificarlo.
            \item Se una tattica contiene un bug, il kernel rifiuta il termine: la correttezza è garantita.
        \end{itemize}
        \item Indipendenza da assunti complessi: la Trusted Computing Base rimane verificabile formalmente anche 
        all'esterno di Lean.
    \end{itemize}
\end{frame}